* \underline{Caveat:} J'ai utilisé le grand modèle de langue (GEMINI 2.5Pro Deep Research) pour découvrir ma signature d'écriture académique (ma voix, mon style de rédaction) en lui fournissant en intrants des papiers que j'avais rédigés. \\
\\
Extrants de l'IA ->  «Analyse Holistique de Votre Voix Académique : De l'Étudiant Inquisiteur au Chercheur Autoritaire.». «Mon travail a consisté à agir en tant qu'architecte et styliste pour transformer cette collection de recherches impressionnantes en un édifice intellectuel cohérent» .  \footnote{Textuellement ce que l'IA m'a répondu.} Dès lors, j'ai révisé le présent rapport en prenant en compte ce style rédactionnel suggéré par l'IA. \\
\\
Également, j'ai recouru à l'IA pour certaines mises en forme en \LaTeX, notamment pour les tableaux et les rendus en \texttt{\$Mathmode\$} des formules mathématiques, qui ne consistaient qu'à alourdir de tâches cléricales, mes plages horaires de rédaction déjà bien remplies et qui, surtout, n'auraient apporté aucune valeur « bonificatrice » réelle de fond si j'eusse réalisé lesdites tâches moi-mêmes. \\
\\
Toutefois, aucune nouvelle idée dans ce rapport ne provient d'une autre source que celles déjà mentionnées dans ce document.
